\documentclass[11pt]{article}
\usepackage{amsmath,amssymb}
\usepackage[margin=1in]{geometry}
\allowdisplaybreaks
\usepackage[UTF8]{ctex}
\begin{document}
\section*{具有时间排序服务功能的移动包裹储物柜，以及可重新加载的仓库设施}
这是\emph{“使用移动包裹柜进行最后一公里配送的路线规划”模型的正式批准版本。该模型采用了混合Q学习网络方法}。它保留了论文中的位置选择、任务构建以及三个成本组成部分。其打印版路线规划约束条件的修改内容如下所述。原始文本则单独存档。

\subsection*{停车位及服务任务}
每位接受服务的客户$n$在每次停靠点$k$提供坐标$h_{nk}$和时间窗口$[E_{nk},L_{nk}]$。通过停车中心$g_i$和标签$b_{nk}$，位置阶段（1)--(3）将停靠点分配给最近的停车中心，以最小化集群内的平方欧几里得距离；每个被占用的停车中心是其分配停靠点的算术平均值。论文选择了最小数量的$P$个站点，这些站点允许步行距离$d_{in}^k\le\rho_n$和服务半径$\rho_n\le r_m$。在该数量的站点中，平方距离目标决定了集群的划分。对于被释放的家族，每位客户有一个停靠点；所有点都是共线的，两个半径都是1，且严格分开的组有唯一的最近停车中心。每组都有一个代表，他们之间的距离超过两个距离单位，从而为$P$提供了一个匹配的下限。路线阶段使用这些停车中心，并不会为了降低交付成本而重新定位它们。

对于每一个允许停车间隔为$[E_i,L_i]$的停车位$i$，间隔范围（15–18）会与停车窗口和客户窗口相交，并将有效可用时间划分为若干个任务，每个任务的持续时间为$SV_i$。在这些情况下，每个客户窗口都包含了一个相关的停车间隔，其宽度等于一个服务周期，因此每个地点有一个任务。不同的任务$u=(i,a)$具有特定的位置$i(u)=i$、服务周期$s_u=SV_i$、时间窗口端点$e_u,l_u$，以及给定的包裹数量$q_u$。邻近性和总需求参数是输入参数；本文并未提供仅根据客户坐标来估算包裹数量的通用公式。此外，不收取任何停车费或占用费。

\subsection*{修订后的路由决策和可行集合}
让$M$代表候选储物柜，$\Psi$代表最大活跃舰队规模，$Q$代表公共储物柜的容量，而$0$则代表仓库。$F_m$表示储物柜$m$的固定成本。对于任何给定的运输路线$(i,j)$，距离为$d_{ij}$，单位旅行成本为$c_{ij}$，而速度为$v>0$。定义$\tau_{ij}=d_{ij}/v$为某个储物柜的部署状态。一个已部署的储物柜$m$会接收到一个有限的序列$R_m=(r_{m0},\ldots,r_{mH_m})$。该储物柜的起始和结束位置都是仓库。每个内部操作要么是服务任务，要么是返回以重新加载物品的过程。连续两个仓库之间的空转行程被排除在外。未使用的储物柜具有$x_m=0$的值，并且没有路线。每个已使用的储物柜具有$x_m=1$的值，并且至少有一个任务。每个任务会在所有路线中只出现一次。
\[
 \sum_m x_m\le\Psi,\qquad
 \sum_m\sum_{k:r_{mk}=u}1=1\quad(u\in U).
\]
对于每一轮$S$，在两次连续的仓库出现之间，需要完成$\sum_{u\in S}q_u\le Q$项任务。后续轮次的包裹可以在储物柜到达仓库时立即装载。在释放的输入中，没有额外的装载时间。让$i(0)=0$、$s_0=0$……在时间零点时，储物柜是可用的。服务开始的时间点$t_{mk}$（以及特定仓库的时间点）是连续且按时间顺序进行的。
\[
 t_{m0}=0,\qquad
 t_{m,k+1}\ge t_{mk}+s_{r_{mk}}+\tau_{i(r_{mk}),i(r_{m,k+1})}
 \quad(0\le k<H_m).
\]
Slack表示等待时间。重新加载是一个按时间顺序进行的访问过程，每个访问都有自己到达/离开的时间点。对于序列$u,0,w$，最早的开始时间是在$w$，这包括两段行程以及$u$处的服务。提前或延迟的服务都是可能的，但会被扣分；任务的时间窗口是灵活的。任务数据由原始的客户和停车窗口决定，而不是由额外的硬性路线时间限制来决定。

\subsection*{目标}
使用$(z)_+=\max(z,0)$时，可以最小化
\[
 \begin{aligned}
 &W_1\sum_m F_m x_m+
 W_2\sum_{m:x_m=1}\sum_{k=0}^{H_m-1}
       d_{i(r_{mk}),i(r_{m,k+1})}\,c_{i(r_{mk}),i(r_{m,k+1})}\\
 &\qquad{}+ W_3\sum_m\sum_{k:r_{mk}\in U}
       \left[(e_{r_{mk}}-t_{mk})_++(t_{mk}-l_{r_{mk}})_+\right].
 \end{aligned}
\]
每个实例都保留其自己的$W_1,W_2,W_3$、单位、费率、舰队以及客户数量。路线构建过程中，会计算所有中途停留和最终返回的时间；固定成本则只支付一次，即每次激活存储设备时。专门指定的数量和服务启动延迟仅适用于实际访问情况。不需要进行零惩罚的旅行。

\subsection*{与印刷版的方程式和作者算法相比，有以下差异：}
物理页面16中给出了(4)--(14)。在(5)中，索引后的固定成本被打印在总和之外。这里的成本是针对每个被激活的储物柜单独计算的。方程(6)是通过对每次实际移动进行收费来实现的。方程(7)则是对所有没有分配系数的储物柜进行求和；这里只有实际进行的访问才会产生影响。方程(8)和(10)的意图保持不变。方程(9)限制了储物柜的总分配数量；而批准的重新加载模型则限制了每次从一个仓库到另一个仓库的行程次数。在(11)--(12)中，自由/总和储物柜索引被替换为每个活跃储物柜的实际开始和结束时间，以及与之对应的中间仓库访问时间。方程(13)中的“所有对都相等”条件也限制了未使用的路径，并且当一个仓库的访问需要不同的时间时，会使用不同的仓库标签。这些都被实际的优先顺序所取代。路径序列则编码了(14)中的二进制访问/顺序/部署选择。

关于物理 p.~13的操作过程的描述，包括初始加载以及一轮运行后的返回过程。第21页到第22页讨论了存储和重新加载的过程；附录A中，第~49页对任务顺序进行了解析。这些段落为修订提供了依据。不过，这些段落并不使修订后的模型与(4)--(14)完全相同，也无法证明该模型的运行结果确实达到了全局最优解。修订后的时钟在之前的存储过程返回之前无法重新启动。历史数据结果则单独附在后面。
\end{document}
